\section{Infrastructure d'exécution et circuit de livraison}
\label{sec:infrastructure}

\subsection{Topologie du laboratoire}

L'exécution de LANCE repose sur un hôte Proxmox unique et un conteneur LXC qui
porte toute la plateforme. Le playbook
\code{benchmarks/ansible/playbooks/deploy\_master.yml} crée, sur le groupe
d'hôtes Proxmox, un conteneur LXC non privilégié d'identifiant 200 nommé
\code{nato-master} (Debian 13, 20~Go, 4096~Mo de mémoire, 2~cœurs,
\code{nesting=1}, \code{onboot 1}). Deux interfaces sont attachées :
\code{eth0} sur \code{vmbr0} en DHCP pour le management, et \code{eth1} sur
\code{vmbr1} en \code{192.168.100.200/24} pour le réseau de test des scénarios.
Le playbook ajoute l'autorisation de périphérique TUN et le montage
\code{/dev/net/tun}, requis par Tailscale, puis reconstruit une chaîne
\code{OUTPUT} d'\code{iptables} : sur \code{eth0} seuls sont autorisés
l'interface locale, les états \code{ESTABLISHED,RELATED}, le DNS, HTTP, HTTPS,
SSH, l'API Proxmox (port 8006), UDP 41641 et le cache apt ; tout le reste est
rejeté. Sur \code{eth1}, la sortie est intégralement autorisée.
\code{eth1} n'a pas de passerelle, ce qui confine les scans au réseau isolé.

Le nom d'hôte Proxmox est donné comme \code{pve-nato} par \code{AGENTS.md},
\code{objectifs/ETAT-COURANT.md} et \code{docs/guides/umons-llm.md}, mais ni le
playbook ni \code{docs/guides/development-instances.md} ne le nomment :
l'inventaire Ansible versionné est un gabarit dont l'adresse vaut littéralement
\code{<PROXMOX\_IP>}, et la seule constante lisible dans le dépôt est
l'identifiant de conteneur, 200. La configuration réelle vit hors du dépôt.

\subsection{Les trois instances}

Le script \code{scripts/deployment/deploy-instance.sh} fixe le triplet
emplacement, port, répertoire et service, et refuse tout autre nom. Le service
stable est produit par le playbook ; les unités de développement sont écrites
par le script.

\begin{table}[H]
\centering
\caption{Instances déployées sur le conteneur maître.}
\label{tab:instances}
\begin{tabular}{@{}llll@{}}
\toprule
Emplacement & Branche & Port & Checkout et service \\
\midrule
\code{main} & \code{main} & 8501 & \code{/opt/nato-smartcity-iot} / \code{nato-fastapi} \\
\code{dev-1} & \code{dev/1} & 8502 & \code{/opt/lance-dev-1} / \code{lance-dev-1} \\
\code{dev-2} & \code{dev/2} & 8503 & \code{/opt/lance-dev-2} / \code{lance-dev-2} \\
\bottomrule
\end{tabular}
\end{table}

Chaque emplacement possède son environnement Python, sa base, son fichier
\code{.env} et ses artefacts. La configuration partagée est écrite dans
\code{/var/lib/lance/<slot>.env} et importée par une surcharge systemd ; ces
entrées alimentent \code{GET /api/environment}, qui renvoie l'emplacement, la
branche, le commit et l'activation du verrou. Le script vérifie après démarrage
que la réponse HTTP annonce exactement le commit validé et l'emplacement attendu,
et que le verrou est actif.

Le verrou commun est un fichier unique, \code{/var/lib/lance/lab.lock}, ouvert
par \code{exec 9>>} puis verrouillé par \code{flock -x 9} avant toute mutation ;
le commentaire est explicite : ne jamais supprimer ce fichier. La même valeur est
transmise aux API, ce qui permet aux trois instances de se coordonner. L'attente
est coopérative, sans ordre d'arrivée, et conservée en mémoire. Un déploiement
d'un emplacement de développement vérifie que l'instance stable annonce le
verrou : une instance stable antérieure n'y participerait pas.

\subsection{Le runner et le circuit de livraison}

Le runner GitHub Actions est installé dans \code{/opt/actions-runner},
enregistré avec le libellé \code{nato-master} et exécuté comme service systemd.

Le circuit se lit en quatre maillons. Un \code{push} sur \code{main},
\code{dev/1} ou \code{dev/2}, ou une demande de fusion vers ces branches,
déclenche \code{.github/workflows/benchmark-integrity.yml}. Le premier job
exécute \code{.github/scripts/ci\_scope.py}, qui classe les fichiers modifiés et
publie le drapeau \code{full}. Le second construit l'image du worker isolé et
vérifie l'absence de vérité terrain. Les jobs de déploiement appellent ensuite
\code{.github/workflows/update-master.yml}, qui lance le script de déploiement
avec l'emplacement, le SHA validé, la branche et le dépôt distant.

\begin{table}[H]
\centering
\caption{Déclencheurs et parcours du circuit.}
\label{tab:circuit}
\small
\begin{tabular}{@{}p{3.1cm}p{5.6cm}p{4.8cm}@{}}
\toprule
Événement & Parcours & Déploiement \\
\midrule
\code{push} documentaire & classement à \code{full=false} : diff du commit précédent et contrôle d'espaces & aucun \\
\code{push} de code & tests complets, contrôles de scénario et Ansible, sandbox, image du worker & \code{main} vers 8501, \code{dev/*} vers 8502 ou 8503 \\
Demande de fusion & même parcours selon classement & aucun \\
Déclenchement manuel & \code{full=true} forcé & emplacement choisi, ou contrôles seuls \\
Étiquette de version & \code{full=true} forcé & \code{main} uniquement \\
\bottomrule
\end{tabular}
\end{table}

Le parcours léger couvre les fichiers Markdown, LaTeX, PDF et images sous
\code{docs/}, \code{objectifs/} et \code{research/}, ainsi qu'une liste fermée
de fichiers de présentation. Le classement conserve les suppressions et les deux
chemins d'un renommage ; tout événement indéterminé bascule sur le parcours
complet, et un échec bloque le déploiement. Un push documentaire ne redémarre
donc aucune instance.

Côté machine, le script de mise à jour récupère le SHA et applique une fusion en
avance rapide pour \code{main}, ou un détachement après contrôle de l'arbre
propre pour les branches de développement ; il vérifie ensuite que
\code{HEAD} égale le SHA validé. Avant cela, le service est arrêté, le verrou
repris, et l'instance interrogée pour refuser de muter un travail actif ou en
attente. En cas d'échec après l'arrêt, le service reste arrêté : aucun retour
automatique à une ancienne version n'est effectué.

\subsection{Corrections du script de déploiement}

L'historique du script montre une séquence serrée de commits le 29 octobre
2026. La version fondatrice vérifiait l'identité de l'hôte par une adresse
codée en dur et exigeait l'utilisateur root.

\begin{table}[H]
\centering
\caption{Corrections successives et cause traitée.}
\label{tab:fixes}
\small
\begin{tabular}{@{}p{1.6cm}p{3.4cm}p{7.4cm}@{}}
\toprule
Commit & Correction & Cause traitée \\
\midrule
\code{6b0a9ff} & adresse en variable d'environnement, ré-exécution \code{sudo}, contrôles de l'instance stable et de la configuration devenus optionnels, CLI Codex sous \code{/usr/local/bin} & un conteneur reconstruit cassait tous les déploiements : identité, privilèges et instances voisines supposés hérités \\
\code{433e2a3} & identité déduite de l'interface \code{tailscale0}, à défaut de variable & supprime la dernière constante dépendante d'une reconstruction \\
\code{a568c68} à \code{dcd5903} & chemin absolu, exécution via \code{bash}, transport du nom de runner et de tout l'environnement de déploiement & \code{sudo} ne résout pas un chemin relatif ni ne conserve l'environnement \\
\code{150eb8a} & échec de téléchargement du CLI Codex toléré & sous \code{set -e}, l'échec arrêtait le déploiement et laissait l'instance arrêtée \\
\bottomrule
\end{tabular}
\end{table}

Chaque correction correspond à un mode de défaillance concret, non à un
raffinement esthétique. La même logique d'identité s'applique désormais aux
contrôleurs de benchmark, avec une variable de suspension explicite, et un test
dédié la couvre.

\subsection{Le modèle UMONS et son relais}

Le fournisseur \code{ollama-umons} est atteint par un transfert local SSH sur
l'hôte Proxmox, vers le cluster universitaire
(\code{docs/guides/umons-llm.md}). Le trajet a été identifié et rétabli le
4~octobre 2026, avec une génération minimale réussie ; l'authentification
observée est interactive. La commande de rétablissement est documentée, ainsi
que les contrôles d'écoute. Le guide est explicite : ni démarrage au boot, ni
reconnexion après la fin du processus ; l'administration du serveur GPU en aval
reste à confirmer avec l'administrateur. La fragilité est constatée dans les
faits : le relais est retombé le 22 octobre, et c'est le pré-vol fournisseur qui
a nommé la panne sans exécuter de bras.

\subsection{La panne du 29 octobre}

Les pannes réellement tracées dans le dépôt pour cette journée sont celles de la
reconstruction du conteneur, qui rendit caduques l'adresse figée et l'hypothèse
d'un runner privilégié. En revanche, \emph{aucune mention d'une route par défaut
multipath sur le routeur du laboratoire n'existe dans le dépôt} : la recherche
sur les termes correspondants ne renvoie rien. Ce fait, constaté en séance mais
non consigné, ne doit pas être présenté comme documenté. Il constitue une
lacune de la mémoire du projet, puisque la même panne se reproduira au prochain
redémarrage du routeur.

\subsection{Fragilités connues}

\begin{itemize}
\item \textbf{Points de défaillance non répliqués} : hôte Proxmox unique,
conteneur unique, runner unique. L'enregistrement du runner tolérant l'échec, un
runner absent ne bloque pas le playbook mais bloque tout déploiement.
\item \textbf{Relais UMONS} : aucun service supervisé, authentification
interactive, adresse encore figée dans les exemples du guide alors que le
commit \code{e54da24} affirmait avoir cessé de figer les adresses dans la
documentation. Redémarrer LANCE ne répare pas le relais.
\item \textbf{Verrou coopératif} : sans ordre d'arrivée, en mémoire, non couvert
pour un script Ansible lancé directement, un ancien checkout ou une intervention
manuelle. Ce n'est pas une isolation de sécurité entre collaborateurs.
\item \textbf{Reprise manuelle} : un échec après arrêt laisse le service arrêté,
sans retour automatique.
\item \textbf{Identité de l'hôte} : sans variable déclarée, l'identité dépend de
la présence d'une adresse quelconque sur l'interface Tailscale ; le script ne
vérifie pas que cet hôte est bien celui qui doit héberger l'instance.
\item \textbf{Documentation partielle} : le nom de l'hôte Proxmox et le segment
physique ne correspondent qu'imparfaitement au playbook ; l'inventaire versionné
est un gabarit.
\end{itemize}
